%----------------------------------------------------------------------------------------
%  A1 portrait poster, visual version: Implementing Conditional Access on FPGA
%  Diagram-first, large type. Compile with pdfLaTeX (Overleaf); figures in figures/,
%  logo.png next to this file. Figures are drawn by figures/make_figs_visual.py
%  (new diagrams) and ../poster/figures/make_figs.py (uaf_steps).
%  If it spills onto page 2, lower \figscale.
%----------------------------------------------------------------------------------------

\documentclass{article}
\usepackage[paperwidth=594mm,paperheight=841mm,margin=10mm]{geometry}

\usepackage[T1]{fontenc}
\usepackage{newtxtext,newtxmath}   % Times
\usepackage{anyfontsize}

\usepackage[dvipsnames]{xcolor}
\definecolor{navy}{HTML}{1F3A5F}
\definecolor{rulegray}{HTML}{B5B5B5}

\usepackage{multicol}
\setlength{\columnsep}{12mm}
\setlength{\columnseprule}{0.6pt}
\renewcommand{\columnseprulecolor}{\color{rulegray}}

\usepackage{graphicx}
\graphicspath{{figures/}}

\pagestyle{empty}
\setlength{\parindent}{0pt}
\setlength{\parskip}{6pt}

% Type scale: large, for reading from a distance
\newcommand{\bodysize}{\fontsize{24}{30}\selectfont}
\newcommand{\headsize}{\fontsize{42}{48}\selectfont}

\newcommand{\figscale}{1.0}

\newcounter{psec}
\newcommand{\psection}[1]{%
  \par\vspace{8pt}\stepcounter{psec}%
  {\headsize\bfseries\color{navy}\thepsec\hspace{0.45em}#1\par}%
  \vspace{-4pt}{\color{navy}\rule{\columnwidth}{1.6pt}}\par\vspace{2pt}}

% a figure without a caption number: #1 width, #2 file
\newcommand{\fig}[2][1.0]{\par\vspace{3pt}{\centering\includegraphics[width=#1\linewidth]{#2}\par}\vspace{3pt}}

\begin{document}
\bodysize

%----------------------------------------------------------------------------------------
%  HEADER
%----------------------------------------------------------------------------------------
\noindent
\begin{minipage}[c]{0.12\linewidth}
  \IfFileExists{logo.png}{\def\haslogo{1}}{\IfFileExists{figures/logo.png}{\def\haslogo{1}}{}}%
  \ifdefined\haslogo
    \includegraphics[width=\linewidth,height=66mm,keepaspectratio]{logo.png}%
  \else
    \fbox{\parbox[c][60mm][c]{0.9\linewidth}{\centering [LOGO]}}%
  \fi
\end{minipage}%
\begin{minipage}[c]{0.76\linewidth}\centering
  {\fontsize{64}{72}\selectfont\bfseries\color{navy} Implementing Conditional Access on FPGA\par}
  \vspace{4mm}
  {\fontsize{30}{36}\selectfont A dual-core, cache-coherent RISC-V processor that frees memory immediately\par}
  \vspace{3mm}
  {\fontsize{30}{36}\selectfont Aryan Kumar, Nishant Dhankhar, Kudumula Vinod Kumar Reddy\qquad\textit{Advisor:} Ajay Singh\par}
  \vspace{3mm}
  {\fontsize{26}{32}\selectfont\itshape Department of Computer Science and Engineering, IIT Gandhinagar\par}
\end{minipage}%
\begin{minipage}[c]{0.12\linewidth}\hfill\end{minipage}

\vspace{3mm}
{\color{navy}\rule{\linewidth}{2.5pt}}
\vspace{1mm}

\begin{multicols}{2}
\bodysize

%----------------------------------------------------------------------------------------
\psection{The problem: use-after-free}
Lock-free data structures let threads read nodes that other threads are removing.
\fig[\figscale]{uaf_steps.png}

%----------------------------------------------------------------------------------------
\psection{Today's fixes delay freeing}
\fig[\figscale]{approaches.png}

%----------------------------------------------------------------------------------------
\psection{Conditional Access (CA)}
New instructions~[1] that reuse the cache-coherence protocol: a write by another core
\textbf{revokes} the reader's access, so memory can be freed \textbf{immediately}.
\fig[\figscale]{instructions.png}
\fig[\figscale]{revoke_flow.png}

%----------------------------------------------------------------------------------------
\psection{Why build it on an FPGA?}
CA was only simulated. Real hardware answers what a simulator cannot:
the \textbf{latency} of the instructions, the \textbf{cost of the coherence events} CA relies on,
and the \textbf{area} it needs. Measured on the KR260:
\fig[\figscale]{measured.png}

\columnbreak

%----------------------------------------------------------------------------------------
\psection{What we built on CVA6~[2]}
\fig[0.76]{architecture_big.png}

%----------------------------------------------------------------------------------------
\psection{Challenges}
\fig[0.84]{challenges.png}

%----------------------------------------------------------------------------------------
\psection{Verification}
\fig[0.84]{verification.png}

%----------------------------------------------------------------------------------------
\psection{Results on the board}
\fig[0.9]{headline.png}
\fig[0.68]{throughput_big.png}

%----------------------------------------------------------------------------------------
\psection{Conclusion}
CA works on real hardware: immediate freeing, fastest under contention, 4.5\% of a core.
\textbf{Next:} a branch-on-failure instruction.

{\fontsize{15}{18.5}\selectfont
[1] A.~Singh, T.~Brown, M.~Spear. Efficient hardware primitives for immediate memory reclamation
in optimistic data structures. arXiv:2302.12958, 2023.
[2] F.~Zaruba, L.~Benini. The cost of application-class processing. \emph{IEEE TVLSI}, 2019.\par
\textbf{Acknowledgements.} We thank Prof.\ Ajay Singh for his guidance throughout this project, and the
Department of Computer Science and Engineering, IIT Gandhinagar, for access to the AMD KR260 board and tools.\par}

\end{multicols}
\end{document}
